% ============================================================
%  Chapter 6 — Results
% ============================================================
\section{Results}
\label{ch:results}

This chapter presents the numerical outcomes of the six scenarios defined in
\cref{ch:scenarios}.  Results are reported at the \emph{loop level}: one
complete execution of the Baltic Sea service cycle (all arcs traversed once).
Annual figures --- relevant for regulatory reporting --- are obtained by
multiplying loop values by the annual loop frequency of approximately 17
loops per year, consistent with the EU MRV observation period.

\subsection{Overview: Scenario Comparison}
\label{sec:res:overview}

\Cref{tab:scenario_comparison} summarises the key performance indicators
across all scenarios.  Attained \CII{}, total voyage cost, and fuel
consumption are reported per loop; compliance slack is defined as
\begin{equation}
  \text{slack} = 1 - \frac{\ciiatt}{\ciithr},
  \label{eq:slack}
\end{equation}
so that a positive value indicates headroom below the threshold.

\begin{table}[h]
\centering
\caption{Scenario comparison --- per-loop results}
\label{tab:scenario_comparison}
\begin{tabular}{llrlrrr}
  \toprule
  ID & Name
     & {\CII{} (g/\GT{}$\cdot$nm)}
     & Rating
     & {Slack (\%)}
     & {Total cost (EUR)}
     & {Fuel (t)} \\
  \midrule
  S1  & As-is validation   & 11.2909 & B & ---   & 743{,}613 & 975.9 \\
  S2  & Max speed          & 14.0354 & C & ---   & 924{,}370 & 1{,}213.1 \\
  S3  & CII-optimal 2026   & 10.2705 & A & 33.9  & 820{,}709 & 887.7 \\
  S4  & Strict 2030 + ETS  & \multicolumn{5}{l}{\textbf{INFEASIBLE} --- min.\ achievable \CII{} = 10.3187~\gGTnm{}
       (target 9.9291~\gGTnm{})} \\
  S4b & Relaxed 2030 + ETS & 10.2705 & C & 14.4  & 1{,}087{,}758 & 887.7 \\
  S5  & Flexible schedule  & 8.7116  & A & 43.9  & 696{,}140 & 752.9 \\
  \bottomrule
\end{tabular}
\end{table}

S3 and S4b cost figures include EU~\ETS{} charges (EUR~40/t\,CO$_2$ in 2026;
EUR~100/t\,CO$_2$ in 2030).  S1 and S2 are evaluated at EUR~800/t fuel with no
carbon price for comparability.

\subsection{S1: As-Is Validation}
\label{sec:res:s1}

S1 replays the mean observed speed on each arc and serves as the empirical
benchmark.  The attained \CII{} of 11.2909~\gGTnm{} corresponds to a
Rating~B, consistent with the actual 2022 attained \CII{} of
11.430253~\gGTnm{} recorded in the EU MRV dataset.  The 1.2\,\% discrepancy
between S1 and the reported value arises from three sources: (i)~rounding of
per-arc speeds to the nearest integer knot; (ii)~use of per-arc mean cargo
weights rather than the full voyage-level cargo distribution; and
(iii)~exclusion of port-side auxiliary consumption, which the fuel model does
not capture.

The EU MRV dataset records a total CO$_2$ emission of
16{,}342.18~t for the observation period (L159--L160 of the EU MRV Voyages
sheet: 15{,}727.27~t HFO + 614.91~t MGO).  S1 produces
3{,}038.9~t per loop, corresponding to an annualised figure of
approximately 51{,}660~t at 17~loops per year --- broadly consistent with
the multi-year MRV record when adjusted for the partial-year observation
window.  The per-loop fuel consumption of 975.9~t and voyage cost of
EUR~743{,}613 provide the economic baseline against which the optimised
scenarios are compared.
Given these modelling limitations, the 1.2\,\% error at the
\CII{} rate level is considered satisfactory validation of the fuel model.

\paragraph{Key insight.}
The as-is operation already achieves a Rating~B \CII{}, suggesting that
the vessel is not significantly over-emitting relative to 2026 targets
even without explicit optimisation.  The binding constraint under S1 is
not \CII{} compliance but rather schedule adherence.

\subsection{S2: Max Speed Upper Bound}
\label{sec:res:s2}

Forcing 16~kn on every arc yields an attained \CII{} of 14.0354~\gGTnm{},
a Rating~C.  This establishes the worst-case emissions upper bound and
quantifies the maximum regulatory exposure if schedule pressures were to
force full-speed sailing on all legs.

Comparing S2 and S3 quantifies the \emph{cost of \CII{} compliance}: the
fuel and cost savings from reducing speed to meet the 2026 Rating~C
threshold versus operating at maximum throughput.  The difference is
reported in \cref{sec:res:s3}.

\subsection{S3: CII-Optimal 2026}
\label{sec:res:s3}

S3 minimises voyage cost subject to the 2026 Rating~C \CII{} threshold
(15.5412~\gGTnm{}) and the published per-arc time budgets.  The optimiser
returns an attained \CII{} of 10.2705~\gGTnm{} --- a Rating~A --- with
33.9\,\% compliance slack.  Total voyage cost is EUR~820{,}709 per loop,
comprising EUR~710{,}141 in fuel (887.7~t at EUR~800/t) and
EUR~110{,}569 in EU~\ETS{} charges (2{,}764~t CO$_2$ at EUR~40/t;
2026 \ETS{} phase-in at 100\,\% coverage).
The annualised cost at 17~loops per year is EUR~13.95~million.

The large slack reveals a structurally important finding: \textbf{the
per-arc schedule is the binding constraint in S3, not the \CII{} target}.
On four bottleneck arcs (\cref{tab:bottleneck_arcs}), the time budget is so
tight that the minimum feasible speed already yields a \CII{} well below the
2026 threshold.  Relaxing the \CII{} constraint further would not permit any
speed increase on these arcs; conversely, tightening it (towards Rating~B)
would require impossible sub-minimum speeds on those arcs.

\paragraph{Practical implication.}
Under the 2026 regulatory framework and the published timetable, the operator
faces essentially no trade-off between cost and \CII{} compliance: the
schedule-optimal solution is simultaneously \CII{}-compliant by a wide
margin.  Investment in \CII{} monitoring or speed management tooling yields
limited additional benefit until the post-2026 reduction pathway becomes
binding.

\subsection{S4: Strict 2030 + ETS (Infeasible)}
\label{sec:res:s4}

S4 demonstrates that a 2030 Rating~B \CII{} target (9.9291~\gGTnm{}) is
\textbf{structurally infeasible} under the published schedule.  The minimum
achievable attained \CII{} --- computed analytically via
\texttt{\_compute\_min\_achievable\_cii()} by running each arc at the lowest
speed consistent with its time budget --- is 10.3187~\gGTnm{}, which exceeds
the Rating~B threshold by 0.39~\gGTnm{} (3.9\,\%).

The infeasibility originates from four bottleneck arcs
(\cref{tab:bottleneck_arcs}) whose published time budgets force minimum
speeds of 14--15~kn, generating disproportionately high fuel consumption per
nautical mile.  Even if all non-bottleneck arcs were operated at
10~kn, the aggregate \CII{} cannot be reduced to the Rating~B level.

\begin{table}[h]
\centering
\caption{Bottleneck arcs forcing minimum speeds in S4}
\label{tab:bottleneck_arcs_s4}
\begin{tabular}{llrrr}
  \toprule
  Arc (i $\to$ j) & Description & {Budget (h)} & {Min.\ speed (kn)} & {\CII{} contrib.} \\
  \midrule
  SEMMA $\to$ SEHLD & Mälarhamn $\to$ Helsingborg & 46.1 & 15 & 0.648 \\
  SEMMA $\to$ SEUME & Mälarhamn $\to$ Umeå        & 46.0 & 15 & 0.646 \\
  SEIGG $\to$ DEKEL & Iggesund $\to$ Kiel         & 47.1 & 14 & 0.519 \\
  SETUN $\to$ NLRTM & Sundsvall $\to$ Rotterdam   & 69.8 & 15 & 0.981 \\
  \bottomrule
\end{tabular}
\end{table}

CII contribution (g/\GT{}$\cdot$nm) is the arc's CO$_2$ mass divided by
the effective vessel capacity and arc distance, evaluated at the minimum
feasible speed.  Bottleneck arcs contribute disproportionately because high
minimum speeds generate cubic fuel penalties on long legs.

\paragraph{Thesis contribution.}
The identification and analytical proof of this structural infeasibility ---
that no speed assignment satisfying the published schedule can achieve
Rating~B by 2030 --- constitutes an original finding of this thesis.  It
demonstrates that \CII{} compliance for stringent future targets depends not
only on vessel technology or fuel choice but on \emph{route and schedule
redesign}.

\subsection{S4b: Relaxed 2030 + ETS}
\label{sec:res:s4b}

S4b retains the 2030 economic environment (EU \ETS{} at EUR~100/t\,CO$_2$,
fuel price EUR~914/t) but relaxes the target to Rating~C
(12.0488~\gGTnm{}).  The optimiser finds a feasible solution with attained
\CII{} = 10.2705~\gGTnm{} and 14.4\,\% slack.

The cost uplift relative to S3 reflects both the higher 2030 fuel price and
the inclusion of EU \ETS{} charges.  The speed profile is identical to S3
(schedule is still the binding constraint), but the cost per loop is
substantially higher due to carbon pricing.

\paragraph{Carbon cost exposure.}
The speed profile in S4b is identical to S3 (the schedule remains the binding
constraint), but the cost increases substantially relative to S3.  The fuel
cost rises from EUR~710{,}141 to EUR~811{,}336 per loop owing to the higher
2030 fuel price (EUR~914/t versus EUR~800/t).  The EU~\ETS{} surcharge at
EUR~100/t\,CO$_2$ adds EUR~276{,}422 per loop (EUR~4.70~million per year),
bringing the total voyage cost to EUR~1{,}087{,}758 per loop
(EUR~18.49~million per year).  This compares to EUR~820{,}709 per loop under
S3, representing a 32.5\,\% cost increase attributable jointly to fuel price
escalation and carbon pricing.

\subsection{S5: Flexible Schedule}
\label{sec:res:s5}

S5 replaces per-arc time budgets with a single global cycle-time cap of
$T_{\max} = 1{,}240$~h and is otherwise identical to S3.  The attained
\CII{} of 8.7116~\gGTnm{} achieves a Rating~A with 43.9\,\% slack.

The key outcome is that fuel consumption and cost are reduced relative to S3:
by eliminating the need to rush on bottleneck arcs, the vessel can slow down
on those legs and redistribute the time savings to other arcs.  The
difference between S5 and S3 quantifies the \emph{value of schedule
flexibility}.

\paragraph{Value of schedule flexibility.}
Comparing S5 and S3 quantifies the value of schedule flexibility directly.
S5 reduces per-loop fuel consumption by 134.7~t (from 887.7~t to 752.9~t),
a saving of 15.2\,\%.  Total voyage cost falls from EUR~820{,}709 to
EUR~696{,}140 per loop (EUR~124{,}569 savings, 15.2\,\%).  Annually, this
represents EUR~2.12~million in avoided fuel and carbon costs.
The attained \CII{} improves from 10.2705 to 8.7116~\gGTnm{}, a reduction
of 1.5589~\gGTnm{} (15.2\,\%), achieved purely through schedule
redistribution with no technology investment.

The 43.9\,\% compliance slack in S5 versus 33.9\,\% in S3 confirms that
schedule flexibility not only reduces costs but creates substantial
regulatory headroom for the tightening post-2026 pathway.  An operator
who can negotiate relaxed port windows would arrive at 2030 with a more
comfortable CII trajectory.

\subsection{Annual Scaling}
\label{sec:res:annual}

Operating at 17~loops per year, the annual regulatory and economic figures
are obtained by scaling the per-loop values in \cref{tab:scenario_comparison}.
\Cref{tab:annual_summary} presents the annualised totals.

\begin{table}[h]
\centering
\caption{Annualised scenario outcomes (17 loops per year)}
\label{tab:annual_summary}
\begin{tabular}{llrrr}
  \toprule
  ID & Rating & {Annual cost (EUR)} & {Annual fuel (t)} & {Annual CO$_2$ (t)} \\
  \midrule
  S1  & B & 12{,}641{,}422 & 16{,}590 & 51{,}661 \\
  S2  & C & 15{,}714{,}283 & 20{,}622 & 64{,}218 \\
  S3  & A & 13{,}952{,}061 & 15{,}091 & 46{,}992 \\
  S4  & \multicolumn{4}{l}{INFEASIBLE} \\
  S4b & C & 18{,}491{,}881 & 15{,}091 & 46{,}992 \\
  S5  & A & 11{,}834{,}373 & 12{,}800 & 39{,}859 \\
  \bottomrule
\end{tabular}
\end{table}

Annual CO$_2$ is computed as annual fuel (t) $\times$ 3.114~kg\,CO$_2$/kg\,HFO.
S1 and S2 annual cost figures assume EUR~800/t fuel and zero carbon price.
S4b annual cost includes both the higher fuel price (EUR~914/t) and
EU~\ETS{} at EUR~100/t\,CO$_2$.

Annual figures are the appropriate metric for EU \ETS{} cost exposure
(which is levied on annual emissions) and for the \CII{} rating (which
is assessed annually).  Per-loop figures are the appropriate metric for
operational decisions at the voyage-planning level.
